\Needspace*{29\baselineskip}
\section{Default Separation and Pointer Aliasing}
\label{app:qcp-memory}

QCP needs initial memory predicates that justify dereferences. The assertion \texttt{Require emp} describes an empty heap, so symbolic execution of the swap below stops at the first read of \texttt{*px}: it cannot derive the precondition of that memory read. A single specification can cover both pointer layouts by stating them as alternatives:

\begin{tcolorbox}[casebox,title={Swap with aliased or disjoint pointers},before skip=9pt,after skip=10pt]
\begin{lstlisting}[style=caseinnerstyle,basicstyle=\ttfamily\fontsize{7.5}{8.5}\selectfont]
void swap(int *px, int *py)
/*@ With x y
    Require (px == py && x == *px) ||
            (px != py && x == *px && y == *py)
    Ensure  (px == py && x == *px) ||
            (px != py && y == *px && x == *py)
*/
{
    int t = *px;
    *px = *py;
    *py = t;
}
\end{lstlisting}
\end{tcolorbox}

In the alias branch, \texttt{px == py} and one points-to predicate describe a shared cell. In the other branch, QCP translates the two dereferences into points-to predicates joined by separating conjunction ($*$), which describes disjoint cells. Symbolic execution produces a return proof obligation for each branch; these obligations still require proof. Because a Frama-C/WP contract without separation leaves aliasing unconstrained, representing its valid overlapping and disjoint inputs in QCP would require explicit heap alternatives, so the default precondition generator selects only the disjoint branch to limit symbolic-execution branching. An aliasing input can be supported by replacing the default separation precondition with an appropriate aliasing predicate, as in the first branch above.
